\documentclass{article}
\usepackage{amsmath}
\usepackage{xcolor}
\begin{document}

ECUACIONES DIFERENCIALES CON COEFICIENTES CONSTANTES NO HOMOGENEA




\begin{gather*}
a_{n}(x)y^{n}+a_{n-1}(x)y^{n-1}+.....+a_{1}(x)y'+a_{0}(x)y=\text{\colorbox[RGB]{245,166,35}{$g$}}\text{\colorbox[RGB]{245,166,35}{$($}}\text{\colorbox[RGB]{245,166,35}{$x$}}\text{\colorbox[RGB]{245,166,35}{$)$}}
\end{gather*}


1) Se resuelve como una E.D. homogenea


\begin{gather*}
a_{n}(x)y^{n}+a_{n-1}(x)y^{n-1}+.....+a_{1}(x)y'+a_{0}(x)y=0 \\
\text{Aca se determina y}_{1}, y_{2}.....y_{n} \\
\text{de tal manera que encontramos y}_{h}
\end{gather*}


2) Se determina la solucion particular 
\begin{gather*}
y_{h}
\end{gather*}
 con alguno de los 

dos metodos ( Variacion de Parametros o Coeficientes Indeterminados)

3) La solucion final esta dada por : 
\begin{gather*}
y = y_{h}+y_{p}
\end{gather*}


\textbf{METODO DE COEFICIENTES INDETERMINADOS}

En este metodo se asume una \textbf{solucion generica} para ciertos casos especificos,

es decir cuando 
\begin{gather*}
g(x)
\end{gather*}
 es una funcion exponencial 
\begin{gather*}
𝛼e^{\text{\colorbox[RGB]{126,211,33}{$𝛽$}}x}
\end{gather*}
, es un juego de 
\begin{gather*}
\text{sen y cos}
\end{gather*}
 o 

es un polinomio.

La representacion de estas funciones puede estar dada de forma individual o se 

puede tener una combinacion de estas funciones por adiciones o diferencias o factores de multiplicacion



\table{Funcion g(x)}{Funcion generica}{\begin{gather*}
\text{\colorbox[RGB]{248,231,28}{$5$}}e^{\text{\colorbox[RGB]{126,211,33}{$5$}}x}
\end{gather*}
  ;   
\begin{gather*}
\text{\colorbox[RGB]{248,231,28}{$\frac{1}{2}$}}e^{\text{\colorbox[RGB]{126,211,33}{$10$}}x}
\end{gather*}
  $\rightarrow$  
\begin{gather*}
𝛼e^{\text{\colorbox[RGB]{126,211,33}{$𝛽$}}x}
\end{gather*}}{\begin{gather*}
Ae^{𝛽x}
\end{gather*}}{}{}


\begin{gather*}
y''+16y=e^{5x}  \rightarrow  y_{1} =cos4x  ; y_{2}= sen4x \\
y_{p} =Ae^{5x}  \\
y''+16y=0 \\
-16sen4x + 16sen4x=0 \\
Ax^{2}+Bx+C \\
Ax^{4}+Bx^{3}+Cx^{2}+Dx+E
\end{gather*}




\image-container




\begin{gather*}
y''+16y= x^{3} * e^{5x}     \\
y_{p} = (Ex^{3}+Bx^{2}+Cx+D)\text{* Ae}^{5x}  \\
y_{p} = (Ex^{3}Ae^{5x}+Bx^{2}Ae^{5x}+CxAe^{5x}+DAe^{5x})   \\
y_{p} = (Ex^{3}Ae^{5x}+Bx^{2}Ae^{5x}+CxAe^{5x}+DAe^{5x})  \\
y_{p} = E*A = F  ; BA=G \\
y_{p} = (Fx^{3}+Gx^{2}+Hx+I)* e^{5x}  \\
\vspace{0.5em} \\
\end{gather*}


 Ejercicios




\begin{gather*}
y″+4y′+3y=3x \\
\text{determinamos y}_{h} \\
e^{rx}(r^{2}+4r+3) = 0 \\
(r+3)(r+1) = 0 \\
y_{h} = c_{1}e^{-3x}+c_{2}e^{-x} \\
\text{para determinar y}_{p}\text{, planteamos una solucion generica} \\
y_{p}= Ax+B \\
\text{Para deterinar los valores de A y B, el metodo indica que debemos} \\
\text{reemplazar las derivadas y la funcion solucion y}_{p}\text{en la en E.D.} \\
y_{p}' = A \\
y_{p}''=0 \\
0+4A+3(Ax+B)=3x \\
3Ax +(4A+3B)x^{0} =3x +0x^{0} 
\end{gather*}


\image-container




\begin{gather*}
3Ax=3x \\
4A+3B = 0x^{0} \\
A = 1 \\
4A+3B =0 \\
4+3B=0 \\
B =-\frac{4}{3} \\
y_{p} = Ax+B = x-\frac{4}{3} \\
\text{Se realiza la prueba respectiva para verificar que la solucion} \\
\text{encontrada es una solucion particular} \\
y_{p} = x-\frac{4}{3} \\
y_{p}' = 1 \\
y_{p}''=0 \\
\text{reemplazando} \\
0+4(1)+3(x-\frac{4}{3})=3x \\
4+3x-4=3x
\end{gather*}





\begin{gather*}
y″+4y′+3y=2e^{3x} \\
y_{p} = Ae^{3x} \\
y_{p}' = 3Ae^{3x} \\
y_{p}'' = 9Ae^{3x} \\
9Ae^{3x}+4(3Ae^{3x})+3Ae^{3x} =2e^{3x} \\
(9A+12A+3A)\text{\colorbox[RGB]{248,231,28}{$e$}}\text{\colorbox[RGB]{248,231,28}{$^{3x}$}} = 2\text{\colorbox[RGB]{248,231,28}{$e$}}\text{\colorbox[RGB]{248,231,28}{$^{3x}$}} \\
24A=2 \\
A = \frac{1}{12} \\
y_{p} =  \frac{1}{12}e^{3x} \\
y = y_{h}+y_{p} = c_{1}e^{-3x}+c_{2}e^{-x}+ \frac{1}{12}e^{3x}
\end{gather*}







\begin{gather*}
y″−4y′+4y=7sent−\cos{}t. \\
y_{p} = Asent+Bcost \\
\vspace{0.5em} \\
y″−4y′+4y=8sen3t \\
y_{p} = Asen3t+Bcos3t \\
\vspace{0.5em} \\
y″−4y′+4y =x^{2}+5x-3   \rightarrow Ax^{2}+Bx+C
\end{gather*}





\end{document}
